\documentclass[12pt,reqno]{amsart}

\usepackage{amsmath,amssymb,amsthm}
\usepackage{hyperref}
\usepackage{booktabs}
\usepackage{array}

\theoremstyle{plain}
\newtheorem{theorem}{Theorem}[section]
\newtheorem{lemma}[theorem]{Lemma}
\newtheorem{proposition}[theorem]{Proposition}
\newtheorem{corollary}[theorem]{Corollary}
\newtheorem{conjecture}[theorem]{Conjecture}
\theoremstyle{definition}
\newtheorem{definition}[theorem]{Definition}
\newtheorem{example}[theorem]{Example}
\theoremstyle{remark}
\newtheorem{remark}[theorem]{Remark}
\newtheorem{openquestion}[theorem]{Open Problem}

\DeclareMathOperator{\Ham}{Ham}
\DeclareMathOperator{\Aut}{Aut}
\DeclareMathOperator{\tr}{tr}
\DeclareMathOperator{\rank}{rank}

\title{Tournaments: Enumeration, Fourier Analysis, and Path Homology}

\author{Multi-Agent Research Project}

\begin{document}

\begin{abstract}
We report results on tournaments (complete directed graphs) spanning three areas:
(1)~high-performance Burnside enumeration extending 90+ OEIS sequences with
approximately 12{,}000 new terms, including tournaments (A000568, $n=77 \to 200$),
self-complementary graphs (A000171, $n=100 \to 439$), and
$k$-uniform hypergraph families ($k = 3, \ldots, 10$);
(2)~the Walsh--Fourier spectrum of Hamiltonian path counts, giving closed-form
formulas for the Walsh coefficients of $H(T)$ and the transfer matrix $M[a,b]$,
and an elementary proof of the Odd-Cycle Collection Formula
$H(T) = I(\Omega(T), 2)$;
(3)~GLMY path homology of tournaments, proving $\beta_2(T) = 0$ for all tournaments
and discovering $\beta_3 = 2$ at $n = 8$.
We also establish universal congruences for the signed Hamiltonian permanent
and prove that $H(T) \neq 7$ and $H(T) \neq 21$ for all tournaments on all~$n$.
\end{abstract}

\maketitle
\tableofcontents

%% ====================================================================
\section{Introduction}
%% ====================================================================

A \emph{tournament} $T$ on $n$ vertices is a complete directed graph: for each
pair $\{u, v\}$, exactly one of the arcs $u \to v$ or $v \to u$ is present.
There are $2^{\binom{n}{2}}$ labeled tournaments on $n$ vertices.

A \emph{Hamiltonian path} in $T$ is a directed path visiting every vertex
exactly once. We write $H(T)$ for the number of such paths. R\'edei's classical
theorem~\cite{Redei1934} states that $H(T)$ is always odd.

This paper reports results from a large-scale computational and theoretical
investigation of tournaments, organized into four programs:

\begin{enumerate}
\item \textbf{Burnside enumeration.} We develop high-performance algorithms
for counting equivalence classes of combinatorial structures under group actions,
extending 90+ sequences in the OEIS~\cite{OEIS} with approximately 12{,}000 new terms.
The key innovations are LCD-scaled integer accumulation, divisor-signature
M\"obius inversion for $k$-uniform hypergraphs, and generating function methods
for matrix enumeration.

\item \textbf{Walsh--Fourier analysis.} We compute the complete Walsh--Fourier
spectrum of $H(T)$ and the transfer matrix $M[a,b]$, obtaining closed-form
formulas for all Walsh coefficients. This yields an elementary proof of the
Odd-Cycle Collection Formula (OCF), a proof that $M[a,b] = M[b,a]$, and extreme
dimensionality reduction (a 100{,}000$\times$ compression at $n = 7$).

\item \textbf{Path homology.} Using GLMY path homology~\cite{GLMY2012}, we prove
$\beta_2(T) = 0$ for all tournaments---a vanishing result with no known
analogue for other graph families. We discover $\beta_3 = 2$ at $n = 8$,
the first Betti number exceeding~1.

\item \textbf{Signed permanent and spectrum gaps.} We establish universal
congruences for the signed Hamiltonian permanent $S(T)$ and prove that
$H(T) = 7$ and $H(T) = 21$ are impossible for all tournaments on all~$n$.
\end{enumerate}

%% ====================================================================
\section{Burnside Enumeration of OEIS Sequences}\label{sec:enumeration}
%% ====================================================================

\subsection{Background}

Burnside's lemma (or the Cauchy--Frobenius lemma) states that the number of
orbits of a finite group $G$ acting on a set $X$ is
\[
|X/G| = \frac{1}{|G|} \sum_{g \in G} |X^g|,
\]
where $X^g = \{x \in X : g \cdot x = x\}$ is the set of fixed points of $g$.

For counting structures on $n$ elements up to permutation, $G = S_n$ acts on
a set of configurations (tournaments, graphs, hypergraphs, etc.). The standard
approach uses the cycle index: each permutation $\sigma$ is characterized by
its cycle type $\lambda = (1^{a_1} 2^{a_2} \cdots)$, and permutations of the
same cycle type contribute the same fixed-point count. This reduces the sum
from $n!$ terms to $p(n)$ terms (the number of integer partitions of $n$).

\subsection{LCD-scaled integer accumulation}

The Burnside sum typically involves rational arithmetic: each partition $\lambda$
contributes $f(\lambda) / z_\lambda$ where $z_\lambda = \prod_i i^{a_i} a_i!$.
Accumulating these fractions naively requires GCD computations at each step.

Our key optimization: precompute $L = \text{lcm}\{z_\lambda : \lambda \vdash n\}$,
then accumulate $L \cdot f(\lambda) / z_\lambda$ as an integer. The final result
is the accumulated sum divided by $L$.

For tournament enumeration (A000568), we combine this with \emph{bucket
accumulation}: since each partition contributes $c \cdot 2^t$ for some integer $t$,
we maintain separate buckets indexed by $t$ and combine at the end. This keeps
intermediate GMP numbers small, avoiding the size explosion that occurs when
adding numbers of different magnitudes.

\begin{theorem}[A000568 speedup]
LCD-scaled bucket accumulation computes $a(n)$ for tournaments in time
$O(p(n) \cdot n)$ arithmetic operations on numbers of $O(n^2)$ digits, yielding
a $250$--$1600\times$ speedup over standard rational DP at $n = 100$--$200$.
\end{theorem}

\subsection{Divisor-signature M\"obius optimization}

For $k$-uniform hypergraphs, the Burnside sum requires counting orbits of
$k$-element subsets under each cyclic permutation $\sigma^t$. The standard
approach iterates $t = 0, 1, \ldots, L-1$ where $L = \text{lcm}$ of cycle
lengths, costing $O(L)$ per partition.

We replace this with M\"obius inversion on the divisor lattice of $L$. For
each cycle of length $c$ in $\sigma$, only divisors $d \mid c$ with $c/d \leq k$
contribute. This reduces the cost to $O(\tau(L) \cdot k)$ where $\tau$ is the
number-of-divisors function.

\begin{theorem}[Hypergraph speedup]
Divisor-signature M\"obius inversion gives a $64\times$ speedup for $4$-uniform
hypergraphs (A051240) and a $17\times$ speedup for $5$-uniform (A051249),
compared to direct Burnside orbit counting.
\end{theorem}

\subsection{Generating function methods for matrix enumeration}

Counting $n \times n$ matrices with entries from $\{1, \ldots, k\}$ up to
simultaneous row permutation, column permutation, and symbol permutation
requires a sum over \emph{triples} of partitions $(\lambda, \mu, \nu)$ of $n$.
Direct enumeration costs $O(p(n)^3)$, which is prohibitive for $n \geq 25$.

We fix two of the three partitions and use an exponential generating function
to sum over the third. For the pair $(\lambda, \mu)$, define
\[
w_j(\lambda, \mu) = k^{K(\lambda, \mu, j)}
\]
where $K$ counts orbits of the joint action at periodicity $j$. Then the
sum over the third partition $\nu$ is
\[
\sum_{\nu} \frac{|C(\nu)|}{n!} \prod_i w_{c_i}(\lambda, \mu)
= [x^n]\, \exp\!\Bigl(\sum_{j \geq 1} w_j \, \frac{x^j}{j}\Bigr).
\]

\begin{theorem}[Matrix enumeration speedup]
The GF method reduces matrix enumeration from $O(p(n)^3)$ to
$O(p(n)^2 \cdot n^2)$. Combined with OpenMP parallelization ($\sim 7\times$
on 8 cores), this extends A091058 from $n = 15$ to $n = 30+$.
\end{theorem}

\subsection{Summary of extensions}

Table~\ref{tab:oeis} lists the major extensions. Full data and enumerator
source code are available in the repository.

\begin{table}[ht]
\centering
\caption{Selected OEIS sequence extensions.}\label{tab:oeis}
\begin{tabular}{@{}llrrr@{}}
\toprule
Sequence & Description & OEIS had & We computed & New \\
\midrule
A000568 & Tournaments & 77 & 200+ & 123+ \\
A000273 & Digraphs & 65 & 101 & 36 \\
A001174 & Oriented graphs & 50 & 200 & 150 \\
A000595 & Binary relations & 51 & 200 & 149 \\
A000171 & Self-comp graphs & 100 & 439+ & 339+ \\
A051240 & 4-uniform hypergraphs & 19 & 77+ & 58+ \\
A052283 & Digraphs by arcs (triangle) & 2681 & 9020 & 6340 \\
A028657 & $m \times n$ binary matrices & $\sim$1081 & 3000+ & 2000+ \\
A091060 & $n \times n$, 3 symbols, $S_n^3$ & 13 & 53+ & 40+ \\
\bottomrule
\end{tabular}
\end{table}

%% ====================================================================
\section{The Odd-Cycle Collection Formula}\label{sec:ocf}
%% ====================================================================

\subsection{Statement and history}

\begin{definition}
The \emph{odd-cycle conflict graph} $\Omega(T)$ has as vertices all directed
odd cycles of $T$, with two cycles adjacent if they share a vertex. The
\emph{independence polynomial} is $I(G, x) = \sum_{k \geq 0} \alpha_k(G) \, x^k$,
where $\alpha_k$ counts independent sets of size $k$.
\end{definition}

\begin{theorem}[OCF; Grinberg--Stanley~\cite{GrinbergStanley2024}]\label{thm:ocf}
For every tournament $T$, $H(T) = I(\Omega(T), 2)$.
\end{theorem}

This gives $H(T) = 1 + 2\alpha_1 + 4\alpha_2 + \cdots$, immediately implying
R\'edei's theorem. The formula was conjectured computationally and verified
exhaustively through $n = 8$ (fixing one arc and varying the remaining $\binom{8}{2}-1 = 27$, giving $2^{27}$ configurations; combined with relabeling invariance this covers all tournaments). Grinberg and Stanley
proved it using P-partition theory~\cite{GrinbergStanley2024}. We give an
independent proof via Walsh--Fourier analysis in Section~\ref{sec:walsh}.

\subsection{Claim A and Claim B}

\begin{theorem}[Claim B, proved]\label{thm:claimB}
For any tournament $T$ and vertex $v$,
\[
I(\Omega(T), 2) - I(\Omega(T \setminus v), 2) = 2 \sum_{C \ni v} \mu(C),
\]
where the sum is over odd cycles through $v$ and $\mu(C) = I(\Omega(T \setminus v)|_{\mathrm{avoid}\, C \setminus \{v\}}, 2)$.
\end{theorem}

The companion \emph{Claim A}---that the same identity holds with $H(T)$
replacing $I(\Omega(T), 2)$---follows from OCF but was also verified
independently through $n = 8$.

\subsection{Practical speedups from OCF}

The OCF replaces the $O(2^n n^2)$ Held--Karp DP algorithm with trace formulas:
\begin{itemize}
\item $\alpha_1 = t_3 + t_5 + t_7 + \cdots$ (odd-cycle counts, computable in $O(n^3)$
via matrix powers)
\item $\alpha_2$ computable from per-vertex cycle counts via inclusion-exclusion ($O(n^3)$)
\end{itemize}

For moderate $n$, the cycle-counting operations are polynomial-time ($O(n^3)$ via matrix traces) while Held--Karp is exponential ($O(2^n n^2)$), giving significant practical speedups.

\subsection{Spectrum gaps}

\begin{theorem}[THM-079]\label{thm:hgap}
$H(T) \neq 7$ and $H(T) \neq 21$ for any tournament $T$ on any $n \geq 1$.
These are the only permanent gaps below $200$.
\end{theorem}

The proof of $H \neq 21$ uses a ``poisoning graph'' DAG argument: any
cycle structure producing $H = 21$ forces additional cycles that push $H$
above 21. Verified exhaustively at $n = 8$ (all $2^{28}$ tournaments).

%% ====================================================================
\section{Walsh--Fourier Analysis of Tournament Invariants}\label{sec:walsh}
%% ====================================================================

\subsection{Setup}

Encode a labeled tournament $T$ on $n$ vertices as a binary string
$t \in \{0,1\}^m$ where $m = \binom{n}{2}$, recording the orientation of
each edge pair. The Walsh--Fourier transform of a function $f$ on tournaments is
\[
\hat{f}[S] = 2^{-m} \sum_{t \in \{0,1\}^m} (-1)^{t \cdot \chi_S} f(t)
\]
for each subset $S \subseteq [m]$.

\subsection{Complete spectrum of $H(T)$}

\begin{theorem}[THM-071, Walsh diagonalization]\label{thm:walshH}
The Walsh coefficient $\hat{H}[S]$ is nonzero if and only if $S$ is a union
of edge-disjoint even-length paths in $K_n$. In that case,
\[
\hat{H}[S] = \varepsilon \cdot \frac{2^r \cdot (n - 2k)!}{2^{n-1}},
\]
where $|S| = 2k$, $r$ is the number of path components, and
$\varepsilon = \pm 1$ depends on the path orientations.
\end{theorem}

\begin{corollary}[Dimensionality reduction]
At $n = 5$, $H(T)$ has $3$ independent Walsh amplitudes on a $1024$-dimensional
space ($341\times$ compression). At $n = 7$, approximately $20$ amplitudes on a
$2{,}097{,}152$-dimensional space ($\sim 100{,}000\times$ compression).
\end{corollary}

\subsection{Complete spectrum of $M[a,b]$}

\begin{theorem}[THM-080]\label{thm:walshM}
For the transfer matrix entry $M[a,b]$ (Hamiltonian paths from $a$ to $b$):
\[
\widehat{M[a,b]}[S] = (-1)^{\mathrm{asc}(S)} \cdot \frac{2^s \cdot (n-2-|S|)!}{2^{n-2}},
\]
where $s$ is the number of unrooted even-length components of $S$ (components
not containing $a$ or $b$), and the coefficient is nonzero only when
$S \cup \{a,b\}$ forms a valid even-length path union with $|S| \equiv n \pmod{2}$.
\end{theorem}

This formula is manifestly symmetric in $a$ and $b$.

\begin{corollary}[Transfer matrix symmetry, THM-030]
$M[a,b] = M[b,a]$ for all tournaments $T$ and all vertices $a, b$.
\end{corollary}

\subsection{Walsh proof of OCF}

\begin{theorem}[THM-077]\label{thm:walshocf}
The Walsh--Fourier coefficients of $H(T)$ and $I(\Omega(T), 2)$ agree for
every subset $S$. Hence $H(T) = I(\Omega(T), 2)$.
\end{theorem}

\begin{proof}[Proof sketch]
The $H$-side formula (Theorem~\ref{thm:walshH}) gives $\hat{H}[S]$ as a
product of directional signs and a factorial term. For the $I$-side, a
generating function factorization (THM-076) shows $\hat{I}[S]$ decomposes
into cycle-path covering contributions times $E[I(\text{remaining})]$.
An EGF calculation gives $F(t,x) = (1+bt)/(1-at)$, and a hockey-stick
identity shows the total depends only on $m = n - 2k$ and the number of
components $r$, matching the $H$-side exactly.
\end{proof}

This proof is elementary and self-contained, bypassing the P-partition
theory of~\cite{GrinbergStanley2024}.

\subsection{Position Character Decomposition}

\begin{theorem}[THM-068, PCD]\label{thm:pcd}
For each vertex $v$ and Walsh subset $S$ of degree $2k$:
\[
\widehat{M_{\mathrm{pos}}[v]}[S] = (-1)^{[v \in N(S)]} \cdot \frac{\hat{H}[S]}{n - 2k},
\]
where $N(S)$ is the set of odd-offset vertices in the path components of $S$.
\end{theorem}

Proved for all degrees and all odd $n$ via a block-placement argument.
The key insight: all macro-items in the block placement have odd width,
so start parity equals macro-position.

%% ====================================================================
\section{GLMY Path Homology of Tournaments}\label{sec:homology}
%% ====================================================================

\subsection{Background}

The GLMY path homology~\cite{GLMY2012} associates to any directed graph $G$
a chain complex $(\Omega_*, \partial_*)$ where $\Omega_p$ consists of
``allowed $p$-paths'' (sequences of $p+1$ vertices satisfying compatibility
conditions) and $\partial_p$ is the alternating face map. The Betti numbers
$\beta_p = \dim(\ker \partial_p / \operatorname{im} \partial_{p+1})$
are topological invariants of the digraph.

\subsection{The Betti landscape}

Exhaustive computation through $n = 6$ and extensive sampling through $n = 10$
(over 50{,}000 tournaments total) reveals:

\begin{center}
\begin{tabular}{@{}crrrrrr@{}}
\toprule
$n$ & $\beta_0$ & $\beta_1$ & $\beta_2$ & $\beta_3$ & $\beta_4$ & $\beta_5$ \\
\midrule
3 & 1 & 0--1 & 0 & & & \\
4 & 1 & 0--1 & 0 & 0 & & \\
5 & 1 & 0--1 & 0 & 0 & 0 & \\
6 & 1 & 0--1 & 0 & 0--1 & 0 & 0 \\
7 & 1 & 0--1 & 0 & 0--1 & 0--6 & 0 \\
8 & 1 & 0--1 & 0 & 0--2 & 0--5 & 0--1 \\
\bottomrule
\end{tabular}
\end{center}

\subsection{$\beta_2 = 0$: statement and proof}

\begin{theorem}[THM-108/109]\label{thm:beta2}
For every tournament $T$ on $n \geq 3$ vertices, $\beta_2(T) = 0$.
\end{theorem}

\begin{proof}[Proof outline]
By strong induction on $n$. For the induction step, use the long exact
sequence of the pair $(T, T \setminus v)$:
\[
\cdots \to H_2(T \setminus v) \to H_2(T) \to H_2(T, T \setminus v) \to H_1(T \setminus v) \to H_1(T) \to \cdots
\]
By induction, $H_2(T \setminus v) = 0$, so $H_2(T)$ injects into
$H_2(T, T \setminus v)$. It suffices to find a vertex $v$ with
$\beta_1(T \setminus v) \leq \beta_1(T)$ (a ``good'' vertex).

\emph{Case 1:} $\beta_1(T) = 1$. Since $\beta_1 \leq 1$ for all tournaments
(THM-103), every vertex is good.

\emph{Case 2:} $\beta_1(T) = 0$, $T$ not strongly connected. Every vertex
deletion preserves non-strong-connectivity, so $\beta_1(T \setminus v) = 0$
for all $v$.

\emph{Case 3:} $\beta_1(T) = 0$, $T$ strongly connected with a cut vertex $v$.
Then $T \setminus v$ is non-SC, so $v$ is good.

\emph{Case 4:} $\beta_1(T) = 0$, $T$ strongly connected with vertex
connectivity $\geq 2$. An \emph{isolation characterization} (THM-109) shows
that bad vertices are severely constrained: in the all-dominated case, they
have score $0$ or $n-1$, limiting the bad set to at most one vertex. In the
free-cycle case, an adjacency argument gives $n - 5$ good vertices for $n \geq 6$.
Base case $n = 5$ verified exhaustively.
\end{proof}

\begin{remark}
For general directed graphs, $\beta_2 > 0$ is common: 70 out of 59{,}049
oriented graphs on 5 vertices have $\beta_2 > 0$. All such graphs have
\emph{twin vertices} (two vertices with identical neighborhoods), which
tournament completeness forbids.
\end{remark}

\subsection{Supporting results}

\begin{theorem}[THM-103]\label{thm:beta1}
$\beta_1(T) \in \{0, 1\}$ for all tournaments $T$.
\end{theorem}

\begin{theorem}[Rank formula]
$\rank(\partial_2|_{\Omega_2}) = \binom{n}{2} - n + 1 - \beta_1(T)$.
\end{theorem}

\begin{theorem}[Mutual exclusivity, THM-095]
$\beta_1(T) \cdot \beta_3(T) = 0$ for all tournaments $T$ on all $n \geq 3$.
\end{theorem}

\subsection{The $n = 8$ threshold}

At $n = 8$, several patterns break:

\begin{itemize}
\item $\beta_3 = 2$ occurs (0.08\% of tournaments), the first Betti number exceeding 1.
\item $\beta_3 \cdot \beta_4 \neq 0$ can coexist ($\sim 0.15\%$), refuting the
``consecutive seesaw'' hypothesis.
\item The inclusion map $i_*: H_3(T \setminus v) \to H_3(T)$ can have nontrivial
kernel, refuting $i_*$-injectivity.
\end{itemize}

\subsection{Paley tournament homology}

For the Paley tournament $T_p$ ($p$ prime, $p \equiv 3 \pmod{4}$):
\begin{itemize}
\item $\beta_d = p - 1$ at $d = p - 3$, and $\beta_d = 0$ for all other $d \geq 1$.
\item Homotopy type: wedge of $(p-1)$ copies of $S^{p-3}$.
\end{itemize}
Verified: $T_7$ has $\beta = (1, 0, 0, 0, 6, 0)$ and $T_{11}$ has $\beta_8 = 10$.

%% ====================================================================
\section{The Signed Hamiltonian Permanent}\label{sec:signed}
%% ====================================================================

\begin{definition}
Let $B = 2A - J$ be the skew-symmetric signed adjacency matrix (entries $\pm 1$).
The \emph{signed Hamiltonian permanent} is
\[
S(T) = \sum_{\sigma \in S_n} \prod_{i=1}^{n-1} B[\sigma(i), \sigma(i+1)].
\]
\end{definition}

\begin{theorem}[THM-A]
$S(T) = 0$ for all tournaments on even $n$.
\end{theorem}

\begin{theorem}[THM-H, universal congruence]
$S(T) \bmod 2^{n-1}$ depends only on $n$, not on $T$.
\end{theorem}

\begin{theorem}[THM-J, universality criterion]
$S(T)$ is fully independent of $T$ (modulo $2^{n-1}$) if and only if
$s_2(n-3) \leq 1$, where $s_2$ is the binary digit sum. The universal values
are $n = 3, 5, 7, 11, 19, 35, 67, \ldots$
\end{theorem}

\begin{theorem}[THM-I, master identity]
For any set of $k$ non-adjacent positions in a Hamiltonian path:
\[
D_S = \frac{n!}{2^k},
\]
where $D_S$ is the sum over $(2k)!$ orderings of $2k$ vertices at those
positions, weighted by signed edge products.
\end{theorem}

At $n = 5$: $c_0 = S(T)/2^4 = H(T) - 3t_3$, encoding both the path count
and the 3-cycle count in a single quantity.

%% ====================================================================
\section{The Worpitzky/F-Polynomial}\label{sec:fpoly}
%% ====================================================================

For a labeled tournament $T$, define $\mathrm{fwd}(P) = \#\{i : P_i < P_{i+1}\}$
for each Hamiltonian path $P$. The \emph{forward-edge polynomial} is
$F(T, x) = \sum_P x^{\mathrm{fwd}(P)}$.

\begin{theorem}[Complement duality]
$F_k(T) = F_{n-1-k}(T^{\mathrm{op}})$.
\end{theorem}

\begin{theorem}[Mod-2 universality]
$F_k(T) \equiv \binom{n-1}{k} \pmod{2}$ for all tournaments $T$.
\end{theorem}

\begin{theorem}[Moment hierarchy]
$\mathrm{Var}(\mathrm{fwd}) = 2t_3$, and the $2k$-th cumulant of $\mathrm{fwd}$
has leading term $2t_{2k+1}/\binom{n}{2k}$.
\end{theorem}

\begin{conjecture}[Unimodality]
$F(T, x)$ is unimodal for all tournaments $T$. Verified: 50{,}000+
tournaments tested with zero violations.
\end{conjecture}

%% ====================================================================
\section{Open Problems}\label{sec:open}
%% ====================================================================

\begin{openquestion}
Why does $\beta_1(T) = 0$ imply $|\{v : \beta_1(T \setminus v) = 1\}| \leq 3$?
Verified through $n = 10$; no algebraic proof.
\end{openquestion}

\begin{openquestion}
What structural property allows $\beta_3 = 2$ at $n = 8$? What bound replaces
$\beta_3 \leq 1$ for $n \geq 8$?
\end{openquestion}

% Open Question 3 removed: beta_1 * beta_3 = 0 is now PROVED for all n
% (THM-095 conditional on beta_2=0 which is proved by THM-108/109)

\begin{openquestion}
Among all tournaments on $p$ vertices ($p$ prime, $p \equiv 3 \pmod{4}$),
does the Paley tournament $T_p$ maximize $H(T)$?
\end{openquestion}

\begin{openquestion}
Prove unimodality of $F(T, x)$ for all tournaments.
\end{openquestion}

\begin{openquestion}
Find a per-path formula valid at all $n$. The 3-cycle formula works at
$n \leq 5$ but fails at $n = 6$ due to 5-cycle contributions. The correct
generalization incorporating all odd cycle lengths is unknown.
\end{openquestion}

%% ====================================================================
\section{Computational Verification}\label{sec:computation}
%% ====================================================================

All results have been verified computationally:

\begin{center}
\begin{tabular}{@{}llrl@{}}
\toprule
Result & Method & Scale & Status \\
\midrule
OCF: $H = I(\Omega, 2)$ & Exhaustive & $n \leq 8$ ($2^{27}$, one arc fixed) & 0 failures \\
$\beta_2 = 0$ & Exhaustive + sampling & $n \leq 10$ (50k+) & 0 failures \\
Walsh spectrum of $H$ & Exhaustive & $n \leq 6$; sampled $n = 7$ & 0 failures \\
Walsh spectrum of $M$ & Exhaustive & $n \leq 6$; sampled $n = 7$ & 0 failures \\
$H \neq 21$ & Exhaustive & $n \leq 8$ ($2^{28}$) & 0 failures \\
Signed permanent & Exhaustive & $n \leq 7$; sampled $n = 9$ & 0 failures \\
$F$ unimodality & Sampling & $n \leq 8$ (50k+) & 0 failures \\
Paley maximality & Exhaustive at primes & $p = 3, 7, 11$ & confirmed \\
\bottomrule
\end{tabular}
\end{center}

Source code for all computations is publicly available in the project repository.

%% ====================================================================
\begin{thebibliography}{99}
%% ====================================================================

\bibitem{Redei1934}
L.~R\'edei,
\emph{Ein kombinatorischer Satz},
Acta Litt.\ Sci.\ Szeged \textbf{7} (1934), 39--43.

\bibitem{GrinbergStanley2024}
D.~Grinberg and R.\,P.~Stanley,
\emph{Counting Hamiltonian paths in tournaments},
arXiv:2412.10572 (2024).

\bibitem{GLMY2012}
A.~Grigor'yan, Y.~Lin, Y.~Muranov, and S.-T.~Yau,
\emph{Homologies of path complexes and digraphs},
arXiv:1207.2834 (2012).

\bibitem{Moon1968}
J.\,W.~Moon,
\emph{Topics on Tournaments},
Holt, Rinehart and Winston, New York (1968).

\bibitem{Forcade1973}
R.~Forcade,
\emph{Parity of paths and circuits in tournaments},
Discrete Math.\ \textbf{6} (1973), 115--118.

\bibitem{OEIS}
OEIS Foundation Inc.,
\emph{The On-Line Encyclopedia of Integer Sequences},
\url{https://oeis.org} (2024).

\bibitem{ElSahili2020}
A.~El Sahili and M.~Abi Aad,
\emph{Parity of paths and circuits in tournaments},
Discrete Math.\ \textbf{343} (2020).

\bibitem{SchweserStiebitzToft2025}
J.~Schweser, M.~Stiebitz, and B.~Toft,
\emph{The tournament theorem of R\'edei revisited},
arXiv:2510.10659 (2025).

\bibitem{TangYau2026}
X.~Tang and S.-T.~Yau,
\emph{Path homology of circulant digraphs},
arXiv:2602.04140 (2026).

\end{thebibliography}

\end{document}
